% ch12_conclusion.tex: Chapter 12, Conclusion and future work.
% The claim list is Table T20 (tab:T20). Every number is traced in ch12_conclusion_NOTES.md.
\chapter{Conclusion and future work}\label{ch:conclusion}

This work set out to build a model that finds almost all cracks, inactive areas and other important defects in EL images of
solar cells. It ends with a benchmark protocol for EL defect segmentation and one confirmed gain of a recipe on sealed cells of
unseen modules. Table~\ref{tab:T20} lists every claim with its evidence, interval and status, the status in its source's own words; this chapter
reads it by status, states what the evidence does not reach, and names the work that would change the picture.

\section{What the evidence supports}\label{sec:concl-claims}

Of the 22 claims of Table~\ref{tab:T20}, one is confirmed zero-shot on the sealed set, one is confirmed as a few-shot result
inside it (out of fold), one is confirmed on CV5 for known modules only, one passes a selection-biased check, three are measurements
of the data, eight are not detected, two failed, and five are descriptive.

\paragraph{Confirmed.} On E1 (unseen modules, supported, not proven; OP-A only; 101 of 106 cells CSIR), the C2 recipe
(HRNet-W18 U-Net++ with SWAD weight averaging and D4 test-time augmentation), averaged over seeds 5 to 9, beats the full-pool PF
recipe on crack relF1@2 by +0.0430 [+0.0013, +0.1043], with all four guards passing (claim 1). The full cells carry the gain and
on the half-cut cells it is not detected (claim 2); most of the gain is the model (claim 5). The lower end is near 0 and the
smallest effect this read could detect is about 0.074, so the size of the gain is open. On CV5, D4 on the PF baseline is confirmed
(+0.0134 [+0.0066, +0.0207]); on E1 the D4 part is not resolved (claim 6).

\paragraph{Passes, selection-biased.} On CV5 the C2 recipe (CS\_D4) passes its check against PF (+0.0222 [+0.0138, +0.0308];
claim 8), but its readout and views were chosen on the same cells, and the ARTS width guard fails on the 416 clean cells.

\paragraph{Measured.} Strict crack IoU measures the annotators' line-width convention more than the model (claim 17). Crack label
errors exist but are a minority of the image-supported candidates and move no contrast by more than 0.003 (claim 18). Claim 19 holds
on strict crack IoU only: the label release moves PF from 0.320 to 0.365 (+0.045), more than any method gain over PF on CV5 (at most
+0.0333). On relF1@2, the deciding metric, it is not claimed: the release moves PF by +0.025, and the largest CV5 contrast against PF
(+0.0243, printed for information) is 0.97 of that shift.

\paragraph{Not detected or failed.} One checkpoint against C1 (claim 3), the ensemble over its members (claim 4), D4 on the
candidate (claim 7), weight averaging alone (claim 10), inactive on E1 (claim 12), C2 on test50 (claim 13), the effect of the module
leak (claim 15) and training on the older label set (claim 20) are not detected. HRNet alone gains on CV5 but fails its guards (claim
9), and the claim that predicted inactive area predicts power loss better than crack length is not supported (claim 21).

\paragraph{Descriptive.} Two models of the EL papers sit below PF under equal training in 2-seed screens (claim 14), the strict-IoU
gain of weight averaging is a convention effect (claim 11), and the crack relF1@2 of the PF recipe is far lower on unseen sources
and modules than on the pool, as levels with no interval for the gap (claim 16).

\paragraph{The few-shot track.} Inside E1, out of fold over 5 folds by source image (seeds 5 to 7, OP-A, 106 cells and 71
clusters, 101 of them CSIR), the C2 recipe trained with the labelled E1 cells of the other folds beats the same recipe zero-shot
(opening \#8's C2, seeds 5 to 7) on crack relF1@2 by +0.2213 [+0.0977, +0.3504], with both guards passing: CONFIRMED (claim 22).
The half-cut cells carry it (+0.2530 [+0.0878, +0.3885]); on half-cut cells without scuff it is not separated ($-$0.0701
[$-$0.3019, +0.0517]). Half of the labelled images give about the same gain (HALF $-$ ZS +0.2042 [+0.0718, +0.3366]; ALL $-$
HALF not separated). The claim is that labelled cells of a layout help, not that the model handles a new layout without them.

\section{The goal, answered}\label{sec:concl-goal}

The operating target set at the start, relR@2 of at least 0.90 at relP@2 of at least 0.90, is not reached: on E1 the C2 recipe
(mean over seeds 5 to 9, OP-A) reaches relP@2 0.5995 and relR@2 0.4157, and even on CV5, in distribution and selection-biased,
0.9216 and 0.8211. Inactive gains
are not detected on E1 or on CV5, and material and edge dark score 0.0000 on E1 for every arm. Our reading (judgement): the lasting
result of this work is the instrument rather than the model. It shows that published EL segmentation numbers cannot be compared,
measures why (release, split leakage, label width, metric conventions, no intervals), and replaces them with a protocol under which
one recipe gain survived a sealed test. On the half-cut layout the zero-shot gain is not detected, and labelled cells of that
layout raise crack relF1@2 out of fold inside E1 (+0.2530 [+0.0878, +0.3885], a stratum that decides nothing), with half of the
labelled images giving about the same gain; on this evidence, a new layout calls for some labelled cells of it.

\section{Future work}\label{sec:concl-future}

Each item below would change a specific limit of Chapter~\ref{ch:limitations}.
\begin{itemize}
  \item \textbf{A second annotation.} An independent annotation of a few dozen pool cells would turn the one double annotation into
  a measured ceiling, and planted errors would measure the recall of the label-error detector; an expert look of 2 to 3 hours at
  the 49 pool cells with an image-supported candidate (confirm or reject only) would check the AI raters' verdicts with a person.
  \item \textbf{A new unseen-module set.} E1 is spent by the few-shot track, so the next unseen-module claim needs new labelled
  cells from new modules, with half-cut cells among them and the module identity recorded at capture.
  \item \textbf{A tighter instrument.} A split-variance probe would let the margin cover the choice of folds as well as seed noise,
  and HOUSE predictions with D4 and a FLOP-matched PF comparator would give the C2 recipe like-for-like and equal-cost contrasts.
  \item \textbf{The weak classes.} OP-C, the per-class bias layer that raised mIoU-R13 on the pool, or a trained arm aimed at
  gridline, splice, edge dark and material, belongs in a final model before those classes can be called addressed.
  \item \textbf{A deployable licence.} A commercial model needs the C2 recipe retrained without the 50 BMRK cells, with the effect of
  dropping them measured, and an answer from the BenchmarkEL authors on commercial use.
  \item \textbf{A pilot with the industry partner.} STC I-V per module, EL at a known current and drone thermal at 600~W/m$^2$ or more on about
  31 modules (for a true rank correlation of 0.5 to give an interval above 0), with the primary written before measuring, would
  test whether the severity output says anything about power on modules the industry partner handles.
\end{itemize}
The benchmark itself, prepared for release as lists and code (Chapter~\ref{ch:protocol}), lets any of these be read under the same rules as
the results of this thesis.
